\section{Scaling the Baselines: Astronomy v1L}
\label{sec:scaled}

Every baseline conclusion in Section~\ref{sec:baselines} rests on ten
questions per system. To test which of them survive a larger sample, we
minted a second instance, \textbf{astronomy v1L}: same cutoff, same
retrieval and judge settings, but corpora rebuilt from broadened
manifests (12 past-corpus and 18 future-corpus ADS query sets covering
clouds and hazes, atmospheric escape, phase curves, high-resolution
spectroscopy, and JWST; 4{,}040 past and 5{,}754 frozen-window future
records---1.8$\times$ and 2.4$\times$ the v1 corpora) and baselines
scaled to 125 questions each. The review--future-work extractor is the
exception by necessity: it exhausts the supply of explicitly posed
questions in all 4{,}040 pre-cutoff abstracts at 49---itself a finding;
the community's already-written-down questions are a finite resource.
In total v1L evaluates 424 frozen baseline questions plus the ten
evidence-graph questions re-run as a cross-instance anchor, all
judge-only. The v1 instance and its released data are untouched.

\begin{table}[t]
  \centering
  \small
  \caption{Astronomy v1L results (judge-only). Brackets are exact 95\%
  Clopper--Pearson intervals. Eng.\ = future-attention rate; Ans.\ =
  answered; Ref.\ = premise refuted; $s_1$ = mean top-1 retrieval
  similarity; lag = mean years to earliest cited engagement. The
  evidence-graph row is the ten v1 questions re-evaluated on the v1L
  corpus (anchor), not a scaled submission.}
  \label{tab:scaled}
  \begin{tabular}{@{}lrllllr@{}}
    \toprule
    System & $n$ & Eng. & Ans. & Ref. & $s_1$ & lag \\
    \midrule
    \texttt{evidence\_graph} (anchor) & 10 & 80\% [44,97] & 30\% [7,65] & 10\% [0,45] & 0.682 & 3.5 \\
    B2 direct LLM & 125 & 96\% [91,99] & 15\% [9,23] & 0\% [0,3] & 0.722 & 1.9 \\
    B4 citation leaders & 125 & 87\% [80,93] & 26\% [18,34] & 3.2\% [1,8] & 0.632 & 2.5 \\
    B3 review future work & 49 & 92\% [80,98] & 22\% [12,37] & 0\% [0,7] & 0.567 & 2.4 \\
    B1 random claims & 125 & 73\% [64,80] & 11\% [6,18] & 0\% [0,3] & 0.622 & 2.8 \\
    \bottomrule
  \end{tabular}
\end{table}

Table~\ref{tab:scaled} gives the scaled results. Sample size changes
two of Section~\ref{sec:baselines}'s three conclusions and sharpens the
third---which is the point of running the experiment.

\paragraph{Revised: coverage does discriminate at scale.} At $n=10$
every non-random system sat at 90--100\% engagement and we concluded
coverage separates only the floor. At $n=125$ the rates pull apart:
random claims 72.8\% [64,80], citation leaders 87.2\%, direct LLM
96.0\% [91,99]; random vs.\ direct-LLM engagement differs at
$p<10^{-4}$ and random vs.\ citation leaders at $p=0.007$ (two-sided
Fisher). The v1 reading was a small-sample artifact. What survives is
the ceiling: fluent, topical LLM questions still approach saturation,
so coverage separates the bottom and middle of the range while
compressing the top---it remains unusable as a sole metric.

\paragraph{Revised: premise refutation is not unique to the
evidence-graph system.} At $n=10$ no baseline refuted a premise; at
$n=125$ the citation-leader baseline catches four refutations (3.2\%
[1,8]): the subsolar-water conclusion for HD~209458\,b (the same
\citealp{macdonald2017} result behind q\_008), the methane-depleted
atmosphere of K2-18\,b overturned by JWST, TiO in WASP-121\,b
unconfirmed by later data, and systematic bias found in benchmark
ultracool-dwarf retrievals. Asking ``is the famous result right?''\ of
enough famous results does eventually catch the ones that fall---note
its hindsight-biased selection (Section~\ref{sec:baselines}) works in
its favor here, since post-cutoff citation counts are inflated by
exactly the controversies that produce reversals. Three things remain
true. Refutation is the rarest outcome for every system (random
templates: 0/125, upper bound 2.9\%---refutations are not free; a
question must aim at a contestable claim). The evidence-graph
submission's nominal rate stays highest (10\% vs.\ 3.2\%), but $n=10$
cannot establish superiority ($p=0.32$). And the two routes to a
refutation differ qualitatively: prominence-chasing rediscovers that
famous claims attract scrutiny, while the evidence-graph question
specified the decisive test from pre-cutoff evidence tensions
(Section~\ref{sec:results:case}). The scaled data cannot yet separate
those routes quantitatively; a scaled evidence-graph submission could.

\paragraph{Sharpened: there is a nonzero answered floor.} Random
robustness templates get \emph{answered} 11.2\% [6,18] of the
time---the field re-examines even arbitrarily chosen results at a
measurable base rate. The v1 estimate of that floor (0/10) was too
flattering to every other system: an answered rate is meaningful only
against $\sim$11\%, not zero. Citation leaders (25.6\%) clear the floor
($p=0.005$); direct LLM (15.2\%) does not ($p=0.42$).

\paragraph{Persistent: the contamination signature.} The direct-LLM
baseline keeps the highest similarity to future literature at scale
($s_1=0.722$ vs.\ 0.567--0.682 for all others), the broadest engagement
(96\%, multi-source rate 94\%), and the earliest mean engagement (1.9
years---its cited evidence concentrates in 2021--2022, the years
closest to its training distribution). Scaling revises one part of the
v1 reading: it does convert engagement into answers (15.2\% vs.\ 0/10
at $n=10$), but at a rate statistically indistinguishable from random
templates, despite engaging twice as much of the literature. Breadth
without resolution remains the signature.

\paragraph{Cross-instance anchor: metrics are corpus-relative.}
Re-judging the ten v1 questions on the 2.4$\times$ larger corpus flips
four labels in both directions: two questions gain \labelfmt{answered}
(richer evidence pools), one drops to \labelfmt{not\_addressed} (its
engaging papers pushed out of a top-8 that a larger corpus makes more
competitive), and one premise moves \labelfmt{weakened}$\to$
\labelfmt{still\_plausible}. q\_008's \labelfmt{answered} +
\labelfmt{refuted} reproduces. The lesson is structural: rates are
functions of the (corpus, retriever, judge) triple, so rows are
comparable only within an instance---v1 and v1L rows must never be
ranked against each other, and the frozen-instance design exists
precisely to make the triple explicit.

\paragraph{What scaling did and did not change.} The scaled study
strengthens the benchmark's discriminative claims (coverage now
separates three tiers; answered rates have a measurable floor) and
weakens one system-level claim (refutation exclusivity). It does not
change the evidence-graph submission's standing---its rates are
unchanged and its refutation reproduces---but it narrows what that
standing demonstrates: at current sample sizes, the defensible
statement is that structured evidence analysis found a refutation by
specifying its test in advance, not that it finds refutations at a
higher rate than strong heuristics. Settling the rate question requires
scaling the \emph{submission}, not just the baselines---the first item
on the revised roadmap.
